% Claim proof file for row claim_3_3 -- "AcyclicUnderIntervention".
% Auto-generated stub by scaffold/solving_current_row.py at row start. A
% manager agent dedicated to writing the TeX proof fills in the body below.
% The proof must:
%   - Stay close to the lecture notes' proof if one exists (always search
%     the LN first for a `\begin{proof}` block immediately following the
%     claim; copy / lightly edit when present).
%   - Otherwise use the LN paradigm and cite prior claims by their `ref`.
%   - Be self-contained enough that a mathematician can follow it without
%     re-reading the LN.
%   - Have no `sorry`, `omitted`, or "obvious" shortcuts.
%
% Compiles standalone via the `subfiles` package.

\documentclass[main]{subfiles}

\begin{document}
% ref: claim_3_3 (proof, with statement restated for readability)
% title: AcyclicUnderIntervention
%
% Cross-references: see claim_<ref>_statement_<title>.tex in this folder
% for the corresponding statement.

% --- statement (restated) ---------------------------------------------------
\begin{Rem}
    If $G$ is acyclic then also $G_{\doit(W)}$ is acyclic and a topological order for $G$ is also one for $G_{\doit(W)}$.
\end{Rem}

% --- proof ------------------------------------------------------------------
\begin{proof}
The whole argument rests on a single observation: hard intervention can only \emph{shrink} the edge set. Indeed, by def\_3\_10,
\[
  E_{\doit(W)}
  \;=\;
  E \sm \lC v' \tuh w' \,\middle|\, v' \in G,\; w' \in W \rC
  \;\ins\; E.
\]
(The clause $v' \in G$ on the removed set is automatic given the def\_3\_1 invariant $E \ins (J \cup V) \times V$, so it does not affect the inclusion above.)

\medskip\noindent\emph{Part A (acyclicity).}
Suppose, for contradiction, that $G_{\doit(W)}$ is not acyclic. By def\_3\_6 there is a node $v$ of $G_{\doit(W)}$ admitting a non-trivial directed walk to itself:
\[
  v = v_0 \tuh v_1 \tuh \cdots \tuh v_n = v
  \qquad \text{in } G_{\doit(W)}, \quad n \ge 1,
\]
i.e.\ $(v_k, v_{k+1}) \in E_{\doit(W)}$ for every $k = 0, \dots, n-1$. The edge inclusion $E_{\doit(W)} \ins E$ noted above then gives $(v_k, v_{k+1}) \in E$ for every $k$, so the same vertex sequence
\[
  v = v_0 \tuh v_1 \tuh \cdots \tuh v_n = v
  \qquad \text{in } G
\]
is a non-trivial directed walk in $G$ (def\_3\_4). From its first step $(v, v_1) \in E$ together with the def\_3\_1 invariant $E \ins (J \cup V) \times V$ we read off $v \in J \cup V$, i.e.\ $v \in G$ (def\_3\_2). Thus $v$ is a node of $G$ admitting a non-trivial directed walk to itself, contradicting acyclicity of $G$ (def\_3\_6).

\medskip\noindent\emph{Part B (topological order).}
Assume in addition $W \ins J \cup V$ and that $<$ is a topological order of $G$ (def\_3\_8). We verify, in turn, that $<$ is a total order on $J_{\doit(W)} \cup V_{\doit(W)}$ and that it respects the parents-precede-children condition for $G_{\doit(W)}$.

\medskip\noindent\emph{Step 1: the node sets agree.} By def\_3\_10:
\begin{align*}
  J_{\doit(W)} \cup V_{\doit(W)}
    &= (J \cup W) \cup (V \sm W) \\
    &= J \cup \lp W \cup (V \sm W) \rp \\
    &= J \cup (V \cup W) \\
    &= (J \cup V) \cup W \\
    &= J \cup V,
\end{align*}
where the second equality is associativity of $\cup$, the third uses $W \cup (V \sm W) = W \cup V$, the fourth is commutativity and associativity, and the last uses the precondition $W \ins J \cup V$.

\medskip\noindent\emph{Step 2: $<$ is a total order on the new node set.} Irreflexivity, transitivity, and trichotomy of $<$ on $J \cup V$ are part of $<$'s being a topological order of $G$ (def\_3\_8). By Step 1 the set $J_{\doit(W)} \cup V_{\doit(W)}$ equals $J \cup V$, so the same three properties hold on it without further work.

\medskip\noindent\emph{Step 3: parents precede children in $G_{\doit(W)}$.} Let $v, w \in G_{\doit(W)}$ and suppose $v \in \Pa^{G_{\doit(W)}}(w)$. By def\_3\_5 this means $v \in G_{\doit(W)}$ together with $v \tuh w \in G_{\doit(W)}$, i.e.\ $(v, w) \in E_{\doit(W)}$. The edge inclusion $E_{\doit(W)} \ins E$ noted at the top of the proof then gives $(v, w) \in E$, i.e.\ $v \tuh w \in G$. The def\_3\_1 invariant $E \ins (J \cup V) \times V$ also forces $v \in J \cup V$, i.e.\ $v \in G$ (def\_3\_2). Combining $v \in G$ with $v \tuh w \in G$, def\_3\_5 yields $v \in \Pa^G(w)$. Since $<$ is a topological order of $G$, def\_3\_8 then gives $v < w$. Hence $<$ also satisfies the parents-precede-children condition for $G_{\doit(W)}$.

Combining Steps 1--3, $<$ is a topological order of $G_{\doit(W)}$ in the sense of def\_3\_8.
\end{proof}

\end{document}
